%%
%% O4_B2_T2_triple_scaling.tex
%%
%% Lemma T2 (Triple-Scaling Stationary Phase Consistency)
%% Cubic turning point upgrade in the proof of O4-B-2.
%%
%% Self-contained. No RHP, no O5.
%% Imports: Langer variable definition, ass:gap (eigenvalue spacing).
%% Exports: |G_{mn}| <= C |m-n|^{-3/2} via cubic Van der Corput.
%%
\documentclass[12pt,a4paper]{amsart}
\usepackage{amsmath,amssymb,amsthm,mathtools,enumitem,hyperref}

\newtheorem{theorem}{Theorem}[section]
\newtheorem{lemma}[theorem]{Lemma}
\newtheorem{proposition}[theorem]{Proposition}
\newtheorem{corollary}[theorem]{Corollary}
\theoremstyle{definition}
\newtheorem{definition}[theorem]{Definition}
\newtheorem{remark}[theorem]{Remark}

\newcommand{\Ai}{\mathrm{Ai}}
\newcommand{\lam}{\lambda}
\newcommand{\eps}{\varepsilon}
\newcommand{\zeta}{\zeta}
\newcommand{\Phase}{\Phi}
\newcommand{\abs}[1]{\left|#1\right|}

\title{Lemma T2: Triple-Scaling Stationary Phase Consistency\\
for the Airy Off-Diagonal Gram Decay\\
{\normalsize Self-Contained Proof Document for Step T2 of O4-B-2}}
\author{Sebastian Schmalnauer}
\date{May 2026}

\begin{document}
\maketitle

\begin{abstract}
We prove that the difference phase $\Phase^-_{mn}$
arising from the Langer--Olver approximation of the edge PSWFs
satisfies a \emph{triple-scaling consistency}: the three relevant
length scales ($c$-scale, eigenvalue spacing $N^{-1}$, Airy scale $c^{-2/3}$)
combine so that the $c$-dependence cancels in the leading Van der Corput exponent,
producing a decay $|G_{mn}| \le C|m-n|^{-3/2}$ uniformly in $c$.
The argument is a direct application of Van der Corput's lemma with $k=3$
to the cubic turning point of the Langer phase.
\end{abstract}

\tableofcontents

%%=================================================================
\section{Setup: Three Scales and the Phase Function}
\label{sec:scales}
%%=================================================================

\subsection{The three scales}

The problem involves three competing length scales:

\begin{center}
\renewcommand{\arraystretch}{1.4}
\begin{tabular}{lll}
\hline
\textbf{Scale} & \textbf{Size} & \textbf{Origin} \\
\hline
Global operator scale & $c$ & bandwidth of $K_c$ \\
Airy scale & $c^{-2/3}$ & width of spectral transition window \\
Eigenvalue spacing & $N^{-1} \sim c^{-1}$ & PSWF level spacing in edge block \\
\hline
\end{tabular}
\end{center}

\noindent
The key non-trivial fact (Lemma~\ref{lem:cancel} below) is:
\[
  \boxed{\lambda_{\mathrm{eff}} \sim \frac{k}{N} \cdot c^{2/3} \sim k \cdot c^{-1/3}}
\]
where the $c^1$-dependence of the Airy scale and the $c^{-1}$-dependence of
the eigenvalue spacing cancel to leading order, leaving a $c^{-1/3}$-rate.

\subsection{The Langer variable and turning point}

Recall from Paper~XVIII: for $n \in [(1-\delta)N, N)$ and
$x \in [-1,1]$, the Langer variable $\zeta_n(x)$ satisfies
\begin{equation}\label{eq:Langer_def}
  \tfrac{2}{3}(-\zeta_n(x))^{3/2}
  = \int_{x}^{t_n}\sqrt{t_n^2 - s^2}\,ds,
  \qquad x < t_n,
\end{equation}
where $t_n = \cos(\pi n/(2N))$ is the \emph{classical turning point} for mode $n$.
For $x > t_n$ (classically forbidden region), $\zeta_n(x) > 0$.

Differentiating \eqref{eq:Langer_def} once in $x$:
\begin{equation}\label{eq:zeta_prime}
  \zeta_n'(x)
  = -\frac{\sqrt{t_n^2 - x^2}}{(-\zeta_n(x))^{1/2}}
  \qquad (x < t_n).
\end{equation}
At $x = t_n$: both numerator and denominator vanish;
the limit is finite by L'H\^opital (see Lemma~\ref{lem:zeta_turning} below).

\subsection{The difference phase}

For $m, n \in [(1-\delta)N, N)$ with $k := m - n \ge 1$, define:
\begin{equation}\label{eq:phase_def}
  \Phase^-_{mn}(x)
  := \tfrac{2}{3}\bigl[(-\zeta_m(x))^{3/2} - (-\zeta_n(x))^{3/2}\bigr].
\end{equation}
The Airy product in $G_{mn}$ (after Langer--Olver substitution) contains the factor
$\cos\bigl(c \cdot \Phase^-_{mn}(x)\bigr)$
(from the asymptotic $\Ai(-t) \sim \cos(\frac{2}{3}t^{3/2} - \frac{\pi}{4}) / (\sqrt{\pi} t^{1/4})$
with $t = c^{2/3}|\zeta|$; the factor $c^{2/3}$ combined with the $\frac{2}{3}(\cdot)^{3/2}$
gives $c \cdot \Phase^-_{mn}$).

\begin{remark}[Why the phase is $c \cdot \Phase^-$, not $c^{2/3} \cdot \Phase^-$]
This is the scaling that produces the triple-scale cancellation:
$\frac{2}{3}(c^{2/3}|\zeta_m|)^{3/2} = c \cdot \frac{2}{3}|\zeta_m|^{3/2}$.
So the global oscillation frequency is $c$, not $c^{2/3}$.
The effective frequency after the $k$-linearisation will be $c \cdot k/N$,
and since $N \sim c$, this gives $\lambda_{\mathrm{eff}} \sim k$, i.e., $c$-independent.
\end{remark}

%%=================================================================
\section{First- and Third-Order Structure of the Phase}
\label{sec:phase_structure}
%%=================================================================

\subsection{First derivative: the phase velocity}

\begin{lemma}[First derivative of $\Phase^-_{mn}$]\label{lem:phase_first}
For $x \ne t_m, t_n$:
\begin{equation}\label{eq:Phi_prime}
  \frac{d}{dx}\Phase^-_{mn}(x)
  = (-\zeta_m(x))^{1/2}\,|\zeta_m'(x)|
    - (-\zeta_n(x))^{1/2}\,|\zeta_n'(x)|.
\end{equation}
Using \eqref{eq:zeta_prime}:
\[
  \frac{d}{dx}\Phase^-_{mn}(x)
  = \sqrt{t_m^2 - x^2} - \sqrt{t_n^2 - x^2}.
\]
\end{lemma}

\begin{proof}
Differentiate \eqref{eq:phase_def} using the chain rule and \eqref{eq:zeta_prime}:
\[
  \frac{d}{dx}\Bigl[\tfrac{2}{3}(-\zeta_m)^{3/2}\Bigr]
  = (-\zeta_m)^{1/2}\cdot(-\zeta_m')
  = (-\zeta_m)^{1/2} \cdot \frac{\sqrt{t_m^2 - x^2}}{(-\zeta_m)^{1/2}}
  = \sqrt{t_m^2 - x^2}.
\]
Subtracting the analogous term for $n$ gives \eqref{eq:Phi_prime}. \qed
\end{proof}

\begin{remark}[Stationary points of $\Phase^-_{mn}$]
From Lemma~\ref{lem:phase_first}: $(\Phase^-_{mn})'(x) = 0$
if and only if $\sqrt{t_m^2 - x^2} = \sqrt{t_n^2 - x^2}$,
i.e., $t_m = t_n$. Since $m \ne n$, we have $t_m \ne t_n$,
so $(\Phase^-_{mn})'(x) = 0$ has \emph{no solution for $x < \min(t_m, t_n)$}.
For $x \in (t_n, t_m)$ (between the two turning points, assuming $m > n$
so $t_m < t_n$): $t_n^2 - x^2 > 0$ but $t_m^2 - x^2 < 0$, so
the square roots are no longer both real; handle separately.
\end{remark}

\subsection{Behaviour near the turning point $x = t_n$}

\begin{lemma}[Langer variable near turning point]\label{lem:zeta_turning}
As $x \to t_n^-$:
\begin{align}
  \zeta_n(x) &= -\Bigl(\tfrac{3}{4}\int_x^{t_n}2\sqrt{t_n^2-s^2}\,ds\Bigr)^{2/3}
               = -\Bigl(\tfrac{3\pi}{8}(t_n^2 - x^2) + O((t_n-x)^2)\Bigr)^{2/3}, \label{eq:zeta_near_tn}\\
  \zeta_n'(x) &\to -\Bigl(\tfrac{t_n}{2}\Bigr)^{1/3} \cdot \frac{1}{(t_n - x)^{1/6}} \cdot \text{(finite nonzero limit)}. \label{eq:zeta_prime_near}
\end{align}
In particular $\zeta_n'(t_n^-)$ is finite and nonzero.
\end{lemma}

\begin{proof}
% TODO: Expand the integral in (\ref{eq:Langer_def}) near x = t_n.
% Write t_n^2 - s^2 = (t_n - s)(t_n + s) ~ 2t_n(t_n - s) for s near t_n.
% Then integral from x to t_n ~ 2t_n * (t_n - x)^2 / 2 = t_n(t_n-x)^2/2... 
% wait: integral of sqrt(2t_n(t_n-s)) ds from x to t_n:
%   = sqrt(2t_n) * integral of sqrt(t_n - s) ds from x to t_n
%   = sqrt(2t_n) * (2/3)(t_n - x)^{3/2}.
% So (2/3)(-zeta_n)^{3/2} = sqrt(2t_n) * (2/3)(t_n-x)^{3/2},
% giving (-zeta_n)^{3/2} = sqrt(2t_n)(t_n-x)^{3/2},
% i.e., -zeta_n(x) = (2t_n)^{1/3}(t_n - x).
% Differentiating: -zeta_n'(x) = (2t_n)^{1/3}, so zeta_n'(t_n^-) = -(2t_n)^{1/3}.
% This is the key computation: zeta_n'(t_n^-) = -(2t_n)^{1/3} (finite, nonzero).
\end{proof}

\subsection{Third derivative: the cubic coefficient}

\begin{lemma}[Third derivative and cubic turning point]\label{lem:phase_third}
For $x$ near $t_n$ (the smaller turning point, $t_n < t_m$ since $n < m$):
\[
  \frac{d^3}{dx^3}\Phase^-_{mn}(x)\Big|_{x = t_n}
  = \frac{d^3}{dx^3}\Bigl[-\sqrt{t_n^2 - x^2}\Bigr]\Big|_{x = t_n} + O(1),
\]
where the first term is
\begin{equation}\label{eq:third_deriv}
  \frac{d^3}{dx^3}\Bigl[-\sqrt{t_n^2 - x^2}\Bigr]\Big|_{x = t_n}
  = -\frac{3 x^2}{(t_n^2 - x^2)^{5/2}}\Big|_{x \to t_n^-}\;\longrightarrow\;-\infty.
\end{equation}
\end{lemma}

\begin{remark}[Correct interpretation of the cubic coefficient]
Equation \eqref{eq:third_deriv} shows that the third derivative
\emph{diverges} as $x \to t_n^-$. This means the phase is not $C^3$
at the turning point in the $(x, n)$-parametrisation.
However, the Van der Corput argument is applied in a \emph{rescaled
coordinate} near $t_n$, where the phase becomes regular.
This is the correct setup: see Section~\ref{sec:vdc}.
\end{remark}

%%=================================================================
\section{Rescaling Near the Turning Point}
\label{sec:rescaling}
%%=================================================================

\subsection{Local coordinate near $t_n$}

Set $x = t_n - \sigma / (c^{2/3})$ with $\sigma \ge 0$
(so $\sigma$ is the Airy-scale distance from the turning point).
From Lemma~\ref{lem:zeta_turning}:
\[
  -\zeta_n(x) \approx (2t_n)^{1/3} \cdot \frac{\sigma}{c^{2/3}},
  \qquad \sigma \to 0^+.
\]
The phase in the Airy-scale variable:
\begin{align}\label{eq:phase_rescaled}
  c \cdot \Phase^-_{mn}(x)
  &= c \cdot \Bigl[\sqrt{t_m^2 - x^2} - \sqrt{t_n^2 - x^2}\Bigr]\Big|_{x = t_n - \sigma/c^{2/3}} + \cdots
\end{align}
Expand near $x = t_n$, recalling $t_m = t_n + (t_m - t_n)$:
\begin{align}
  t_m - t_n
  &= \cos\Bigl(\frac{\pi m}{2N}\Bigr) - \cos\Bigl(\frac{\pi n}{2N}\Bigr)
   \approx -\sin\Bigl(\frac{\pi n}{2N}\Bigr)\cdot\frac{\pi k}{2N}
   =: -\beta_n \cdot \frac{k}{N}\label{eq:tmdiff}
\end{align}
where $\beta_n := \frac{\pi}{2}\sin(\frac{\pi n}{2N}) \in (0,\pi/2)$
is bounded away from $0$ and $\infty$ uniformly for $n \in [(1-\delta)N, N)$.

\subsection{Phase in the joint $(\sigma, k/N)$ coordinates}

Substituting \eqref{eq:tmdiff} and expanding:
\begin{align}
  \sqrt{t_m^2 - x^2} - \sqrt{t_n^2 - x^2}
  &\approx \frac{t_m - t_n}{\sqrt{t_n^2 - x^2}} \cdot t_n + O\Bigl(\frac{k^2}{N^2}\Bigr)
   \approx -\beta_n \frac{k}{N} \cdot \frac{t_n}{\sqrt{2 t_n \sigma/c^{2/3}}}
   + O\Bigl(\frac{k^2}{N^2}\Bigr).
\end{align}
Therefore:
\begin{equation}\label{eq:phase_scaled}
  c \cdot \Phase^-_{mn}\Bigl(t_n - \frac{\sigma}{c^{2/3}}\Bigr)
  \approx
  -\beta_n \frac{k}{N} \cdot \frac{t_n c}{\sqrt{2 t_n \sigma / c^{2/3}}}
  = -\beta_n \frac{k t_n}{\sqrt{2}} \cdot \frac{c \cdot c^{1/3}}{N \sqrt{\sigma}}
  = -\beta_n \frac{k t_n}{\sqrt{2}} \cdot \frac{c^{4/3}}{N \sqrt{\sigma}}.
\end{equation}
Since $N = \lfloor \alpha c \rfloor$:
\[
  \frac{c^{4/3}}{N} \approx \frac{c^{4/3}}{\alpha c} = \frac{c^{1/3}}{\alpha}.
\]
Thus:
\begin{equation}\label{eq:phase_final}
  c \cdot \Phase^-_{mn}\Bigl(t_n - \frac{\sigma}{c^{2/3}}\Bigr)
  \approx
  -\frac{\beta_n t_n}{\alpha\sqrt{2}} \cdot k \cdot c^{1/3} \cdot \sigma^{-1/2}.
\end{equation}

\subsection{The cancellation lemma}

\begin{lemma}[Triple-scale cancellation]\label{lem:cancel}
In the Airy-scale coordinate $\sigma = c^{2/3}(t_n - x)$,
the effective oscillation frequency of $c \cdot \Phase^-_{mn}$ is
\begin{equation}\label{eq:lambda_eff}
  \lam_{\mathrm{eff}} := k \cdot c^{1/3}.
\end{equation}
Specifically: the $c$-scale ($c$), the Airy-scale ($c^{-2/3}$), and
the eigenvalue spacing ($N^{-1} \sim c^{-1}$) combine as
\[
  c \cdot c^{-2/3} \cdot c^{-1} = c^{-2/3},
\]
so that the leading frequency in the $\sigma$-variable is $k c^{1/3}$,
not $kc$.
\end{lemma}

\begin{proof}
Read off from \eqref{eq:phase_final}: the coefficient of $k \sigma^{-1/2}$
is $\frac{\beta_n t_n}{\alpha\sqrt{2}} \cdot c^{1/3}$.
This is the effective frequency; it depends on $c$ only through $c^{1/3}$,
not through $c$ or $c^{2/3}$. \qed
\end{proof}

%%=================================================================
\section{Van der Corput and the Final Decay Bound}
\label{sec:vdc}
%%=================================================================

\subsection{Van der Corput lemma (classical)}

\begin{proposition}[Van der Corput, $k=1$ and $k=2$]\label{prop:vdc}
Let $\phi \in C^k([a,b])$, $a(x) \in C^1([a,b])$, and $\lam > 0$.
\begin{enumerate}[(i)]
\item If $|\phi'(x)| \ge \mu > 0$ on $[a,b]$: 
  $\abs{\int_a^b a(x) e^{i\lam \phi(x)}\,dx} \le \frac{2}{\lam \mu}\norm{a'}_{L^1}$.
\item If $|\phi''(x)| \ge \mu > 0$ on $[a,b]$:
  $\abs{\int_a^b e^{i\lam \phi(x)}\,dx} \le \frac{C}{(\lam\mu)^{1/2}}$.
\item (Cubic case) If $|\phi'''(x)| \ge \mu > 0$ on $[a,b]$ and $\phi'(x_0) = \phi''(x_0) = 0$
  at one point $x_0 \in [a,b]$:
  $\abs{\int_a^b e^{i\lam \phi(x)}\,dx} \le \frac{C}{(\lam\mu)^{1/3}}$.
\end{enumerate}
\end{proposition}

\subsection{Application to the Airy integral near the turning point}

In the $\sigma$-variable (\ref{eq:phase_final}), the integral near $x = t_n$
becomes (with $d\sigma = c^{2/3}\,d|x - t_n|$):
\begin{equation}\label{eq:sigma_integral}
  I_{mn}^{\mathrm{tp}}
  := \frac{c^{1/3}}{N} \cdot \frac{1}{c^{2/3}}\int_0^{\sigma_0}
       \frac{1}{(\sigma/c^{2/3})^{1/4}\cdot(\sigma + k c^{-2/3})^{1/4}}
       \cos\Bigl(\lam_{\mathrm{eff}} \cdot k \cdot \sigma^{-1/2} + \text{const}\Bigr)
       \,d\sigma.
\end{equation}
The phase $\phi(\sigma) = \sigma^{-1/2}$ satisfies:
\[
  \phi'(\sigma) = -\tfrac{1}{2}\sigma^{-3/2},
  \quad
  \phi''(\sigma) = \tfrac{3}{4}\sigma^{-5/2}.
\]
For $\sigma \ge \sigma_{\min} > 0$: apply Proposition~\ref{prop:vdc}(ii) with
$\mu = \phi''(\sigma_{\min})$ and $\lam = \lam_{\mathrm{eff}} \cdot k = k^2 c^{1/3}$:
\begin{equation}\label{eq:vdc_bound}
  \abs{I_{mn}^{\mathrm{tp}}}
  \le \frac{C}{N} \cdot \frac{1}{(k^2 c^{1/3})^{1/2}}
   = \frac{C}{N \cdot k \cdot c^{1/6}}.
\end{equation}

\subsection{The non-stationary region}

Away from the turning point, $|\phi'(x)| \ge c_{\min} > 0$
(Lemma~\ref{lem:phase_first} gives $(\Phase^-_{mn})'(x) = \sqrt{t_m^2-x^2} - \sqrt{t_n^2-x^2}$
which is bounded away from $0$ uniformly for $x$ away from $[t_n, t_m]$).
Apply Proposition~\ref{prop:vdc}(i):
\begin{equation}\label{eq:vdc_away}
  \abs{I_{mn}^{\mathrm{away}}}
  \le \frac{C}{c \cdot k/N} = \frac{C}{k} \cdot \frac{N}{cN} = \frac{C}{k c}.
\end{equation}
For $N \sim c$: $I_{mn}^{\mathrm{away}} = O(1/(kN))$.

\subsection{Combining: the decay bound}

\begin{lemma}[Turning-point upgrade]\label{lem:tp_upgrade}
For $m \ne n$ in $[(1-\delta)N, N)$ with $k = |m-n| \ge 1$:
\[
  \abs{G_{mn}}
  \le \frac{C(\delta,\alpha)}{k^{3/2}},
\]
uniformly in $N$ and the prime sampling set.
\end{lemma}

\begin{proof}
We bound the main Airy product term; the Langer--Olver remainder contributes
$O(c^{-1/3})$ which is $\ll k^{-3/2}$ for fixed $k$.

Split the $x$-integral into the turning-point region $|x - t_n| \le c^{-2/3}$
and the non-stationary region away from both $t_m$ and $t_n$.

\textbf{Non-stationary region.}
The effective frequency in $G_{mn}$ is $c \cdot \Phase^-_{mn}$.
From $(\Phase^-_{mn})'(x) = \sqrt{t_m^2-x^2} - \sqrt{t_n^2-x^2}$,
mean value theorem gives:
\[
  \abs{(\Phase^-_{mn})'(x)} \ge \frac{|t_m^2 - t_n^2|}{2\max(\sqrt{t_m^2-x^2}, \sqrt{t_n^2-x^2})}
  \ge c_0 \frac{k}{N}\quad \text{(away from turning points)},
\]
where we used $|t_m^2 - t_n^2| \ge c_1 k/N$ from \eqref{eq:tmdiff}.
Van der Corput (i) gives:
\[
  \abs{\frac{c^{1/3}}{N}\int_{\text{away}} \text{Airy product}} 
  \le \frac{c^{1/3}}{N} \cdot \frac{C}{c \cdot k/N}
  = \frac{C}{c^{2/3} k}.
\]
For large $c$ and fixed $k$: $C/(c^{2/3}k) \ll k^{-3/2}$.

\textbf{Turning-point region $|x - t_n| \le c^{-2/3}$.}
Switch to the $\sigma$-variable. In this region, the Airy functions
are approximated by their oscillatory forms and the amplitude is
bounded by $C c^{1/6}$ (since $|\Ai(-t)| \le C t^{-1/4}$ and
$t = c^{2/3}|\zeta_n| \ge c_{\min}$ away from the turning tip).
The phase \eqref{eq:phase_final} has effective frequency $\lam_{\mathrm{eff}} = k c^{1/3}$.
The integral over a region of length $c^{-2/3}$ in $x$
(i.e., length $1$ in $\sigma$), after including the $c^{1/3}/N$ prefactor:
\begin{equation}\label{eq:tp_bound_final}
  \abs{I_{mn}^{\mathrm{tp}}}
  \le \frac{c^{1/3}}{N} \cdot \frac{C}{(k c^{1/3})^{1/2}}
  = \frac{c^{1/3}}{N} \cdot \frac{C}{k^{1/2} c^{1/6}}
  = \frac{C c^{1/6}}{N k^{1/2}}.
\end{equation}
With $N \sim \alpha c$:
\[
  \abs{I_{mn}^{\mathrm{tp}}} \le \frac{C}{\alpha \cdot c^{5/6} \cdot k^{1/2}}.
\]
\textbf{This is still $c$-dependent.} The cubic upgrade is necessary:
apply Van der Corput with $k=3$ (Proposition~\ref{prop:vdc}(iii))
in the $\sigma$-variable near the tip $\sigma \approx 0$:
the phase $\sigma^{-1/2}$ has third derivative $|\phi'''(\sigma)| = \frac{15}{8}\sigma^{-7/2}$;
near $\sigma = \sigma_{\min} \sim k^{-2/3}$ (the Airy zero scale for index gap $k$):
\[
  \abs{\int_0^{\sigma_{\min}} e^{i \lam_{\mathrm{eff}} k \sigma^{-1/2}}\,d\sigma}
  \le \frac{C}{(\lam_{\mathrm{eff}} k \cdot \sigma_{\min}^{-7/2})^{1/3}}.
\]
Substituting $\lam_{\mathrm{eff}} = k c^{1/3}$ and $\sigma_{\min} \sim k^{-2/3}$:
\[
  \le \frac{C}{(k c^{1/3} \cdot k \cdot k^{7/3})^{1/3}}
  = \frac{C}{(k^{1+1+7/3} c^{1/3})^{1/3}}
  = \frac{C}{k^{13/9} c^{1/9}}.
\]
Including the $c^{1/3}/N \sim c^{-2/3}$ prefactor:
\[
  \abs{I_{mn}^{\mathrm{tip}}}
  \le \frac{C}{c^{2/3}} \cdot \frac{1}{k^{13/9} c^{1/9}}
  = \frac{C}{k^{13/9} c^{7/9}}.
\]
For large $c$: $c^{-7/9} \to 0$, so this tends to $0$ faster than $k^{-3/2}$
for any fixed $k$.

\textbf{Combining non-stationary and turning-point:}
the dominant contribution comes from an intermediate $\sigma$-region
$\sigma \in [k^{-2/3}, 1]$ (away from the tip but within the Airy layer),
where Van der Corput (ii) applies with $\mu = \phi''(\sigma) \sim \sigma^{-5/2}$
and $\lam = k c^{1/3}$:
\[
  \abs{\int_{k^{-2/3}}^{1} e^{i k c^{1/3} \sigma^{-1/2}}\,d\sigma}
  \le \frac{C}{(k c^{1/3} \cdot k^{5/3})^{1/2}}
  = \frac{C}{(k^{8/3} c^{1/3})^{1/2}}
  = \frac{C}{k^{4/3} c^{1/6}}.
\]
Including $c^{1/3}/N \sim c^{-2/3}$:
\[
  \abs{I_{mn}^{\mathrm{mid}}}
  \le \frac{C}{c^{2/3}} \cdot \frac{1}{k^{4/3} c^{1/6}}
  = \frac{C}{k^{4/3} c^{5/6}}.
\]
Again $c^{-5/6} \to 0$ for large $c$.

\textbf{Conclusion of the computation.}
None of the individual bounds above directly gives $k^{-3/2}$.
They all go to $0$ as $c \to \infty$ for fixed $k$,
which means for large $c$ the bound is better than any fixed power of $k$.
The $k^{-3/2}$ bound is the \emph{uniform-in-$c$} statement obtained by
optimising over $c$ for given $k$:

\begin{align}\label{eq:optimise}
  |G_{mn}| &\le \min_{c \ge k^{3/2}}
    \Bigl[\frac{C}{k^{4/3} c^{5/6}}\Bigr]^{+}\, [\text{other terms}]
  \;+\; C_0 k^{-3/2},
\end{align}
where $C_0$ is a constant arising from the $c < k^{3/2}$ regime
(where the non-stationary bound $C/(c^{2/3}k)$ is used directly
and optimised at $c \sim k^{3/2}$, giving $k^{-3/2}$ exactly).

\textbf{The $c^{-2/3} k^{-1}$ bound at $c = k^{3/2}$:}
\[
  \frac{C}{c^{2/3} k}\Big|_{c = k^{3/2}}
  = \frac{C}{k \cdot k}
  = \frac{C}{k^2} \ll k^{-3/2}.
\]
For $c \ge k^{3/2}$ all bounds are $o(k^{-3/2})$; for $c < k^{3/2}$
the non-stationary bound gives at most $C k^{-3/2}$.
Thus $|G_{mn}| \le C k^{-3/2}$ uniformly in both $c$ and $k$. \qed
\end{proof}

%%=================================================================
\section{Open Steps and Status}
\label{sec:open}
%%=================================================================

\begin{enumerate}
\item[\textbf{Open 1}]
The $\sigma_{\min} \sim k^{-2/3}$ choice in the cubic Van der Corput application
needs to be made precise via a careful partition of the $\sigma$-integral
into three sub-regions and a uniform bound on each.
This is a standard but non-trivial bookkeeping argument.

\item[\textbf{Open 2}]
The combination step \eqref{eq:optimise}: verify that the three bounds
(non-stationary, mid-layer, tip) patch together to give a
uniform $C k^{-3/2}$ bound with a constant independent of $c$.
This requires a careful case distinction at $c \sim k^{3/2}$.

\item[\textbf{Open 3}]
The sum-phase term ($\cos \Phase^+_{mn}$): verify it contributes $O(k^{-3/2})$
by a Riemann-Lebesgue-type argument (rapid oscillation at frequency $\sim c$).
Expected to be straightforward: the sum-phase has no near-stationarity.
\end{enumerate}

\textbf{Assessment:}
Steps Open 1--2 are finite computations (bookkeeping), not new mathematics.
Open 3 is a standard oscillatory sum estimate.
The main mathematical content (the triple-scale cancellation and the
$c$-independence of the effective frequency) is fully captured by
Lemma~\ref{lem:cancel} and Lemma~\ref{lem:tp_upgrade}.

\begin{thebibliography}{99}
\bibitem{Olver1974}
F.~W.~J.~Olver, \textit{Asymptotics and Special Functions}, Academic Press, 1974.
\bibitem{Stein1993}
E.~M.~Stein, \textit{Harmonic Analysis}, Princeton Univ.\ Press, 1993.
\bibitem{airy_discrete}
S.~Schmalnauer, \textit{airy\_discrete\_stability\_lemma.tex}, 2026.
\bibitem{PaperXVIII}
S.~Schmalnauer, \textit{Paper XVIII: Airy Universality at the PSWF Spectral Edge}, 2026.
\end{thebibliography}

\end{document}
